%% ============================================================
%%  chapters/01_counseling.tex
%%  Chapter 1: Counseling & Admission Process (COAP / CCMT)
%% ============================================================

\chapter{Counseling \& Admission Process}
\label{chap:counseling}

\begin{protip}[Most Frequently Asked]
  This chapter addresses the most common questions about \COAP (for IITs \& IISc)
  and \CCMT (for NITs). Read it before applying to ensure you don't miss critical
  deadlines or steps.
\end{protip}

%% ─────────────────────────────────────────────────────────
\section{Overview of the Two Portals}
\label{sec:portals}

After the \GATE result is declared, admissions to postgraduate programs are
channeled through two centralized portals depending on the institution:

\begin{tblr}{
    colspec = {X[2,l] X[3,l] X[3,l]},
    hlines,
    vlines,
    row{1} = {bg=pfBlue, fg=white, font=\sffamily\bfseries},
    row{2,4} = {bg=pfGrey},
  }
  \toprule
  Portal & Institutions Covered & Opens Approximately \\
  \midrule
  \textbf{\COAP} (Common Offer Acceptance Portal)
    & IITs, IISc
    & Last week of March \\
  \textbf{\CCMT} (Centralized Counseling for M.Tech)
    & NITs, IIEST Shibpur, select IIITs \& CFTIs
    & Last week of May \\
  \bottomrule
\end{tblr}

\vspace{6pt}
\noindent M.Tech classes typically begin at the \textbf{end of July or August}.

%% ─────────────────────────────────────────────────────────
\section{The COAP Process (IITs \& IISc)}
\label{sec:coap}

\subsection{Step-by-Step Workflow}

\begin{enumerate}[label=\textcolor{pfBlue}{\textbf{Step \arabic*.}}, leftmargin=*, itemsep=6pt]
  \item \textbf{Apply on institute websites.}
        Visit each desired IIT/IISc portal separately, select your preferred
        branch, and submit your application along with your \GATE score.

  \item \textbf{Register on the \COAP portal.}
        Create a single \COAP account---all offer communications will appear here
        regardless of how many institutes you have applied to.

  \item \textbf{Attend interviews (where required).}
        Some institutes conduct interviews before shortlisting. If shortlisted,
        you will receive instructions via email from the respective institute.
        \textbf{Physical attendance is mandatory for interviews.}

  \item \textbf{Check offers each round.}
        Accepted and pending offers become visible after each \COAP round.
        You may \textbf{hold} an offer for up to \textbf{2 rounds} while waiting
        for a better offer.

  \item \textbf{Accept, hold, or decline.}
        From your pool of concurrent offers, select the most preferred one.
        Accepting an offer removes all other offers automatically.

  \item \textbf{Freeze your seat.}
        Once you are satisfied, freeze the offer. The institute will then send
        admission instructions and a fee payment link.
\end{enumerate}

\begin{protip}[Don't lose your hold window]
  An offer can be held for a maximum of \textbf{2 consecutive rounds}. If you
  neither accept nor decline within this window, it is automatically withdrawn.
  Track the round schedule carefully.
\end{protip}

\subsection{Interviews at IITs: Which Institutes Require Them?}

Not all IITs require interviews for M.Tech admissions. The table below summarizes
the current policy (verify on official institute websites each year):

\begin{tblr}{
    colspec  = {X[2,l] X[2,l] X[4,l]},
    hlines, vlines,
    row{1}   = {bg=pfBlue, fg=white, font=\sffamily\bfseries},
    row{2,4,6,8} = {bg=pfGrey},
  }
  \toprule
  Institute & Interview for M.Tech? & Notes \\
  \midrule
  IISc Bangalore
    & \textcolor{pfOrange}{\textbf{Yes (all branches)}}
    & Interview mandatory for every M.Tech \& M.Tech Research branch. \\
  IIT Delhi
    & Partial
    & IEC branch requires interview; check other branches officially. \\
  IIT Kharagpur
    & Partial
    & EE department has interview; ECE department generally does not. \\
  IIT Bombay
    & \textcolor{pfGreen}{\textbf{No (M.Tech)}}
    & No interview for 2-year M.Tech VLSI branches; MS (3-yr) requires interview. \\
  IIT Madras
    & \textcolor{pfGreen}{\textbf{No (M.Tech)}}
    & Similar to IITB; ICS/VLSI M.Tech has no interview. MS and TI-MS require one. \\
  Other IITs
    & Generally No
    & MS/M.Tech Research programs at \textit{all} IITs require interviews. \\
  \bottomrule
\end{tblr}

%% ─────────────────────────────────────────────────────────
\section{The CCMT Process (NITs \& CFTIs)}
\label{sec:ccmt}

The \CCMT process operates similarly to \COAP but serves NITs, IIEST Shibpur,
select IIITs, and Centrally Funded Technical Institutes (CFTIs).

\begin{pfActionList}
  \item Register on the \CCMT portal after \GATE results.
  \item Fill in your choice of institute and branch in order of preference.
  \item Seat allotments happen in multiple rounds; participate in document
        verification and fee payment to confirm your seat.
  \item Check the official \CCMT portal for round-wise cutoffs:
        \url{https://ccmt.admissions.nic.in}
\end{pfActionList}

\begin{protip}[NIT Cutoff Reference]
  A community-sourced NIT cutoff spreadsheet (CCMT 2025) is available at:\\
  \url{https://docs.google.com/spreadsheets/d/1G0KjJN9MNMiwWDRDf4OrzgEBYivGnn\_lT0G7jK57ikk}
\end{protip}

%% ─────────────────────────────────────────────────────────
\section{CFTI Mode Admissions}
\label{sec:cfti}

\begin{eligalert}[CFTI Mode Eligibility]
  IISc and IIT Kanpur (and select other IITs) offer admissions via the
  \textbf{CFTI mode} for certain branches (e.g., Semiconductor Technology at
  IISc). Applicants are typically required to:
  \begin{pfChecklist}
    \item Hold a degree from a CFTI institution.
    \item Have a CGPA of \textbf{8.0 or above}.
    \item Clear an interview (same panel as the GATE-based MS program).
  \end{pfChecklist}
\end{eligalert}

\noindent\textbf{Note on IIT Kanpur VLSI:} It has been observed historically that
all general category seats in the VLSI program are filled by IIT B.Tech graduates,
leaving very limited seats for external \GATE candidates. Verify the latest data
before applying.

%% ─────────────────────────────────────────────────────────
\section{Documents Required}
\label{sec:documents}

Each institute specifies its exact document list in the admission instructions.
The following are universally required across most IITs and NITs:

\begin{pfChecklist}
  \item Class 10 (Secondary) mark sheet \& certificate
  \item Class 12 (Higher Secondary) mark sheet \& certificate
  \item All B.Tech/B.E. semester grade cards
  \item B.Tech/B.E. degree certificate (or provisional)
  \item \GATE scorecard (current year)
  \item Proof of nationality (Aadhaar / Passport)
  \item Category certificate (SC/ST/OBC-NCL/EWS, as applicable)
  \item Recent passport-size photographs
\end{pfChecklist}

%% ─────────────────────────────────────────────────────────
\section{Key Dates \& Timeline}
\label{sec:dates}

\begin{deadline}[GATE 2026 Counseling Timeline]
  \begin{tblr}{
      colspec = {X[3,l] X[3,l]},
      hlines, vlines,
      row{1} = {bg=pfOrange, fg=white, font=\sffamily\bfseries},
      row{2,4,6} = {bg=pfLightOrange},
    }
    \toprule
    Event & Approximate Date \\
    \midrule
    \GATE 2026 Result & March 2026 \\
    \COAP Portal Opens (IITs \& IISc) & Last week of March 2026 \\
    \COAP Counseling Rounds & April -- June 2026 \\
    \CCMT Portal Opens (NITs \& CFTIs) & Last week of May 2026 \\
    M.Tech Classes Begin & July / August 2026 \\
    \bottomrule
  \end{tblr}
\end{deadline}

%% ─────────────────────────────────────────────────────────
\section{Frequently Asked Questions}
\label{sec:counseling-faq}

\begin{faqentry}{If I clear the cutoff for colleges X and Y and apply to both,
  can I choose between them?}
  Yes. In \COAP, you can receive multiple offers from various colleges in any
  round. You are free to choose your preferred branch and college from the pool
  of offers. However, once you \textbf{accept} one offer, all other offers
  disappear. Choose wisely---accept only when you are certain.
\end{faqentry}

\begin{faqentry}{How will I know if I am eligible for an interview at an
  institution?}
  If you are shortlisted, the institution will send instructions directly to your
  registered email address. Apply broadly to all institutes of interest---the
  shortlisting and communication happen automatically.
\end{faqentry}

\begin{faqentry}{Is physical travel required for IIT interviews?}
  \textbf{Yes.} As of current practice, all IIT and IISc interviews require the
  candidate to be physically present at the institute. Online modes have not been
  standardized across all institutes.
\end{faqentry}
